%%%%%%%%%%%%%%%%%%%%%%%%%%%%%%%%%%%%%%%%%%%%%%%%%%%%%%%%%%%%%%%%%%%%%%%%%%%%%%%%
%% Document
%%%%%%%%%%%%%%%%%%%%%%%%%%%%%%%%%%%%%%%%%%%%%%%%%%%%%%%%%%%%%%%%%%%%%%%%%%%%%%
\documentclass[10pt]{beamer}
% latexmk -pdf -pdflatex="lualatex %O %S" enos-slides.tex

% ************************************************************************
% Personal data and user ad-hoc commands
% ************************************************************************
\newcommand{\myName}{%
  Ronan-Alexandre Cherrueau, Adrien Lebre,\newline Dimitri Pertin,
  \underline{Anthony Simonet} and Matthieu
  Simonin\xspace} % 2
\newcommand{\myDepartment}{Inria, France} % 8
\newcommand{\myTime}{\date\xspace} % 9
\newcommand{\myTitle}{Enos: a Holistic Framework for Conducting
  Scientific Evaluations of OpenStack} % 11
\newcommand{\myVersion}{version 0.1\xspace}
\date{May 14th, 2017}
 \titlegraphic{%
	\includegraphics[height=1.0cm]{figs/logo/inria.pdf}\hfill%
	\includegraphics[height=1.2cm]{figs/logo/imta.png}\hfill%
	\includegraphics[height=1.2cm]{figs/logo/ls2n.png}\hfill%
	\includegraphics[height=1.5cm]{figs/discovery-sticker.png}\hfill%
 }

% ********************************************************************
% Setup, finetuning, and useful commands
% ********************************************************************
\newcommand{\ie}{i.\,e.}
\newcommand{\Ie}{I.\,e.}
\newcommand{\eg}{e.\,g.}
\newcommand{\Eg}{E.\,g.}

% ************************************************************************
% Loading some handy packages
% ************************************************************************
%% \PassOptionsToPackage{utf8}{inputenc}
%%  \usepackage{inputenc}
\PassOptionsToPackage{english}{babel}
\usepackage{babel}

% ********************************************************************
% General useful packages
% ********************************************************************
\PassOptionsToPackage{T1}{fontenc} % T2A for cyrillics
	\usepackage{fontenc}
\usepackage{xspace} % to get the spacing after macros right

% ************************************************************************
% code / listings
% ************************************************************************
\usepackage{listings}

\lstdefinelanguage{proverif-pv}{
classoffset=0,
morekeywords={%
  among,channel,choice,clauses,const,def,%
  elimtrue,else,equation,equivalence,event,%
  expand,fail,forall,free,fun,get,if,in,%
  inj-event,insert,let,letfun,new,noninterf,%
  not,nounif,or,otherwise,out,param,phase,%
  pred,proba,process,proof,putbegin,query,%
  reduc,set,suchthat,table,then,type,%
  weaksecret,yield},
otherkeywords={||,&&,->,<->,<=>,==>},
emph={gather},
sensitive=true,
morecomment=[n]{(*}{*)},
morestring=[b]",
%rulesepcolor=\color{solarized@base03},
numbers=left,
classoffset=1,
morekeywords={private,data,typeConverter,%
  memberOptim,decompData,decompDataSelect,%
  block,attacker,mess},
keywordstyle=\bfseries\color{solarized@base00},
classoffset=0
}[keywords,comments,strings]


\lstdefinelanguage{idris}[]{Haskell}{
  morekeywords={total,auto,using,Refl,interface,%
  implementation},
  deletekeywords={intersect,List,Double,Bool,%
  String,nub,fst,snd,not,Either,Maybe,Just,Nothing,%
  either,index,last,id,Eq},
  otherkeywords={=>,>>=},
  morecomment=[l]{|||},
  literate=
    {[]}{{\bot}}2 {->}{→}1 {++}{{+\!\!+}}1
    {<-}{{←}}1
    ,
  mathescape=true,% /!\ takes care of '\' inside a code
}

\lstdefinelanguage{javascript}{
  keywords={typeof, new, true, false, catch, function, return, null,
  catch, switch, var, if, in, while, do, else, case, break},
  ndkeywords={class, export, boolean, throw, implements, import,
  this},
  sensitive=false,
  comment=[l]{//},
  morecomment=[s]{/*}{*/},
  morestring=[b]',
  morestring=[b]"
}

\lstset{
basicstyle=\ttfamily\small,
numbers=none,%left,%
numberstyle=\tiny\color{solarized@base01},
keywordstyle=\bfseries\color{solarized@orange},
stringstyle=\color{solarized@cyan},
%% identifierstyle=\color{solarized@base03},
commentstyle=\color{solarized@base01},
%% emphstyle=\color{solarized@orange},
emphstyle=\bfseries\color{solarized@orange},
fontadjust=true,%
keepspaces=true,%
basewidth={0.58em, 0.53em},%
breaklines=true,
tabsize=2,%
%framerule=0.3pt,%
numbersep=1.5mm,%
upquote=true, % straight quote instead of quotation mark
extendedchars=true,
aboveskip=0pt,
belowskip=0pt,
%% mathescape=true,% /!\ this put a mess in the sh code
%% escapechar=★,%
%% firstnumber=last,
escapeinside={`}{'} % CR as escape char
}


% ************************************************************************
% 6. PDFLaTeX, hyperreferences and citation backreferences
% ************************************************************************
% ********************************************************************
% Using PDFLaTeX
% ********************************************************************
%% I got an error when having a pound `#' symbol in the first argument
%% of an `\href' inside a footnote. The footnote tries to interpret
%% it. The bigfoot package fixes the problem.
%% http://tex.stackexchange.com/a/12860
\PassOptionsToPackage{pdftex,hyperfootnotes=false,pdfpagelabels}{hyperref}
	\usepackage{hyperref} % backref linktocpage pagebackref
\pdfcompresslevel=9
% \pdfadjustspacing=1

% ********************************************************************
% Hyperreferences
% ********************************************************************
\hypersetup{%
    %draft,	% = no hyperlinking at all (useful in b/w printouts)
    colorlinks=false, linktocpage=true, pdfstartpage=3, pdfstartview=FitV,%
    % uncomment the following line if you want to have black links (e.g., for printing)
    %colorlinks=false, linktocpage=false, pdfborder={0 0 0}, pdfstartpage=3, pdfstartview=FitV,%
    breaklinks=true, pdfpagemode=UseNone, pageanchor=true, pdfpagemode=UseOutlines,%
    plainpages=false, bookmarksnumbered, bookmarksopen=true, bookmarksopenlevel=1,%
    hypertexnames=true, pdfhighlight=/O,%nesting=true,%frenchlinks,%
    urlcolor=revealjs1, citecolor=webgreen, %pagecolor=RoyalBlue,%
    %urlcolor=Black, linkcolor=Black, citecolor=Black, %pagecolor=Black,%
    pdftitle={\myTitle},%
    pdfauthor={\textcopyright\ \myName, \myDepartment},%
    pdfsubject={},%
    pdfkeywords={},%
    pdfcreator={pdfLaTeX},%
    pdfproducer={LaTeX with hyperref and classicthesis}%
}


% ********************************************************************
% Last, but not least...
% ********************************************************************
\usepackage{ulem} %to strike out text
\renewenvironment{quote}
{\vspace{.25em}\itshape\list{}{\rightmargin\leftmargin}%
\item\relax\textquotedblleft\ignorespaces}
{\unskip\unskip\textquotedblright\endlist\vspace{.25em}}
% ************************************************************************

% ********************************************************************
% Using different fonts
% ********************************************************************

%% Crimson: https://fontlibrary.org/en/font/crimson
\usepackage[no-math]{fontspec}
%% Scale to math text size
\defaultfontfeatures{Scale=1.05 }
%\newfontfamily\Fonti[
%    Path              = /Users/msimonin/Library/Fonts/,
%    Extension         = .otf,
%    Ligatures         = TeX,
%    UprightFont       = *-Roman,
%    BoldFont          = *-Bold,
%    ItalicFont        = *-Italic,
%    BoldItalicFont    = *-BoldItalic
%]{Crimson}
%\newfontfamily\Fontii[
%    Path              = /Users/msimonin/Library/Fonts/,
%    Extension         = .otf,
%    Ligatures         = TeX,
%    WordSpace         = 1.3,
%    ItalicFont        = Vollkorn-MediumItalic,
%    BoldFont          = Vollkorn-Semibold,
%    BoldItalicFont    = Vollkorn-SemiboldItalic
%]{Vollkorn-Medium}
%\newfontfamily\FontiiD[
%    Path              = /Users/msimonin/Library/Fonts/,
%    Extension         = .otf,
%    Ligatures         = TeX,
%    UprightFont       = *-Regular,
%    BoldFont          = *-Bold,
%    ItalicFont        = *-Italic,
%    BoldItalicFont    = *-BoldItalic
%]{Vollkorn}
%\newfontfamily\Fontiii{Lato}
%
%\setmainfont[
%    Path              = /Users/msimonin/Library/Fonts/,
%    Extension         = .otf,
%    Ligatures         = TeX,
%    UprightFont       = *-Roman,
%    BoldFont          = *-Bold,
%    ItalicFont        = *-Italic,
%    BoldItalicFont    = *-BoldItalic
%]{Crimson}
%
%\setsansfont{Lato}
%
%\setmonofont[
%    Path              = /Users/msimonin/Library/Fonts/,
%    Extension         = .otf,
%    Scale             = MatchLowercase,
%    Ligatures         = TeX,
%    UprightFont       = *-Regular,
%    BoldFont          = *-Medium
%]{FiraMono}%% \linespread{1.05} % a bit more for Palatino


\usepackage[absolute,overlay]{textpos}
\usepackage{enumitem}
\setenumerate[1]{%
  label=\protect\usebeamerfont{enumerate item}%
        \protect\usebeamercolor[fg]{enumerate item}%
        \insertenumlabel.}
\setitemize{label=\usebeamerfont*{itemize item}
  \usebeamercolor[fg]{itemize item}
  \usebeamertemplate{itemize item}}
\setlist{noitemsep,topsep=5pt,parsep=5pt,partopsep=5pt}


\usetheme[pageofpages=of,% String used between the current page and the
                         % total page count.
          titleline=true,% Show a line below the frame title.
          titlepagelogo=figs/logo/all,% Logo for the first page.
          fleuronfamily=EBGaramond, % fleuron font family ❧
          ]{Ronana}

% Set colors
%% \colorlet{colorprimary}{scarletred1}
%% \colorlet{colorsecondary}{scarletred2}
%% \colorlet{colortertiary}{scarletred3}
%% \colorlet{colorquaternary}{colortertiary!50!black}

%% No navigation symbols
\setbeamertemplate{navigation symbols}


%%%%%%%%%%%%%%%%%%%%%%%%%%%%%%%%%%%%%%%%%%%%%%%%%%%%%%%%%%%%%%%%%%%%%%%%%%%
%% Slides
%%%%%%%%%%%%%%%%%%%%%%%%%%%%%%%%%%%%%%%%%%%%%%%%%%%%%%%%%%%%%%%%%%%%%%%%%%%
\begin{document}

% Gray layer after pause
%% \setbeamercovered{transparent}

%%%%%%%%%%%%%%%%%%%%%%%%%%%%%%%%%%%%%%%%%%%%%%%%%%%%%%%%%%%%%%%%%%%%%%%%%
%% Main Page
%%%%%%%%%%%%%%%%%%%%%%%%%%%%%%%%%%%%%%%%%%%%%%%%%%%%%%%%%%%%%%%%%%%%%%%%%
\begin{frame}[t,plain]
  \titlepage
\end{frame}

%%%%%%%%%%%%%%%%%%%%%%%%%%%%%%%%%%%%%%%%%%%%%%%%%%%%%%%%%%%%%%%%%%%%%%%%%
%% The What ?
%%%%%%%%%%%%%%%%%%%%%%%%%%%%%%%%%%%%%%%%%%%%%%%%%%%%%%%%%%%%%%%%%%%%%%%%%

%% Context: number of commits, etc.
%%%%%%%%%%%%%%%%%%%%%%%%%%%%%%%%%%%%%%%%%%%%%%%%%%%%%%%%%%%%%%%%%%%%%%%%%
\begin{frame}{OpenStack}
  \begin{textblock*}{.5\paperwidth}(.2cm,3cm)
  	\centering\textbf{Industry standard} for creating public and private clouds
  \end{textblock*}

  \begin{textblock*}{.6\paperwidth}(.2cm,4.6cm)
   \includegraphics[scale=.65]{figs/dashboard.png}
  \end{textblock*}

  \begin{textblock*}{6cm}(6.5cm,1.8cm)
	 \includegraphics[scale=1.1,angle=350]{figs/os_metrics.png}
  \end{textblock*}

  %\Rightarrow{} every OpenStack is different
\end{frame}

%% Context: services
%%%%%%%%%%%%%%%%%%%%%%%%%%%%%%%%%%%%%%%%%%%%%%%%%%%%%%%%%%%%%%%%%%%%%%%%%
\begin{frame}{OpenStack}
  	\centering\includegraphics[width=.8\paperwidth]{figs/os_big-tent.png}
  \begin{itemize}
	 \item \textbf{164 services} in OpenStack
	 \item Some services are composed of \textbf{sub-services} (e.g. nova-scheduler,
		nova-conductor, \dots)
	 \item Most of which can be disabled/enabled depending on one's needs
  \end{itemize}
\end{frame}

%% Context: topology
%%%%%%%%%%%%%%%%%%%%%%%%%%%%%%%%%%%%%%%%%%%%%%%%%%%%%%%%%%%%%%%%%%%%%%%%%
\begin{frame}{Deployment topologies}
  % 1st figure
  \only<1-3>{
  \begin{textblock*}{.3\paperwidth}(1cm, 1.8cm)
	 \centering
	 \includegraphics[height=3cm]{figs/deployment-scenario1.png}

	 \scriptsize{Single controller, multiple compute nodes}
  \end{textblock*}
  }

  % 2nd figure
  \only<2-3>{
  \begin{textblock*}{.3\paperwidth}(6.5cm, 1.8cm)
	 \centering
	 \includegraphics[height=3cm]{figs/deployment-scenario2.png}

	 \scriptsize{Multiple controllers, multiple compute nodes}
  \end{textblock*}
  }

  % 3rd figure
  \only<3>{
  \begin{textblock*}{.8\paperwidth}(1.0cm, 5.8cm)
	 \centering
	 \includegraphics[height=3.2cm]{figs/deployment-scenario3.png}

	 \scriptsize{Multiple controllers, multiple compute nodes, multiple regions}
  \end{textblock*}
  }
\end{frame}

\begin{frame}{Motivation}

  \textbf{Goals:}
  \begin{itemize}
		\small
	 \item Favour \textbf{reproducible scientific evaluations}
	 \item Automated \emph{Performance Regression Testing}
	 \item Help developers evaluate their code with multiple setups and
		topologies
  \end{itemize}

  \vspace{.3cm}
  \visible<2-3>{
  \textbf{Typical experimental workflow:}
  \begin{enumerate}
		\small
	 \item Book and provision servers
	 \item Deploy OpenStack
	 \item Perform benchmarks
	 \item Collect and save metrics
	 \item Visualize \& share
  \end{enumerate}
  }

  \visible<3>{
	 \Rightarrow{} Automation \& reproducibility means \textbf{all these steps} must be
	 \textbf{scripted}.
  }
\end{frame}


%% Quoi personnelles
%%%%%%%%%%%%%%%%%%%%%%%%%%%%%%%%%%%%%%%%%%%%%%%%%%%%%%%%%%%%%%%%%%%%%%%%%
\begin{frame}{Contribution: Enos}
  \emph{A holistic framework for scientific evaluations of OpenStack}
  \vspace{.4cm}

  \begin{itemize}
	 \item Container based OpenStack deployment with Kolla-ansible-ansible
    \item Rally benchmark (control plane)/Shaker benchmark (data plane)
    \item Metrics collection with cAdvisor and collectd and stored in InfluxDB
    \item Query metrics with Grafana
    \item Query logs with Heka
  \end{itemize}

  Based on \textbf{high level specifications}
\end{frame}

\begin{frame}{Enos: a typical workflow}
\begin{enumerate}
\item \texttt{enos deploy}:
  \begin{itemize}
  \item Gets testbed resources
  \item Deploys OpenStack
  \item Populates OpenStack with cirros image, public/private network
  \end{itemize}
\item \texttt{enos bench}:
  \begin{itemize}
  \item Runs benchmarks
  \item Measures CPU/RAM/Network consumption per site/node/service
  \end{itemize}
\item \texttt{enos backup}:
  \begin{itemize}
  \item Performs post-mortem analysis
  \end{itemize}
\end{enumerate}
\end{frame}


\AtBeginSection[]
{
  \begin{frame}<beamer>{Enos Workflow}
    \tableofcontents[
      currentsection
    ]
  \end{frame}
}

\section{\texttt{enos deploy}}

\begin{frame}[fragile]{\insertsectionhead}
  A \emph{Provider} gets testbed resources
  \begin{itemize}
	 \item Resource: anything running a Docker daemon and Enos can \texttt{ssh} to.
  \item Supports Grid'5000, VirtualBox, OpenStack and Chameleon
  \item Easy to extend: \textasciitilde 500 LOC
  \medskip
  \begin{columns}
    \footnotesize
    \begin{column}{.3\textwidth}
  \begin{lstlisting}
provider: g5k
resources:
  paravance:
    control: 1
    network: 1
    compute: 50
  \end{lstlisting}
    \end{column}
    \begin{column}{.05\textwidth}
      \Rightarrow
    \end{column}
    \begin{column}{.6\textwidth}
      \begin{enumerate}
      \item Get 52 nodes from the paravance cluster;
		\item \texttt{ssh} and install Python and Docker daemon;
      \item Return a list of IP addresses to install OpenStack to.
      \end{enumerate}
    \end{column}
  \end{columns}
  \end{itemize}
\end{frame}

\begin{frame}[fragile]{Topology control: \texttt{reservation.yaml}}
  \begin{textblock*}{.4\paperwidth}(6cm,3cm)
	 \includegraphics[scale=.5]{figs/topology_control.png}
  \end{textblock*}
  \begin{lstlisting}
provider: g5k
resources:
  paravance:
    control: 1
    database: 1
    nova-conductor: 1
    network: 1
    storage: 1
    compute: 50
  \end{lstlisting}
\end{frame}

\begin{frame}[fragile]{Topology control: \texttt{inventory.yaml}}
  \begin{textblock*}{.4\paperwidth}(6cm,3cm)
	 \includegraphics[scale=.5]{figs/topology_control.png}
  \end{textblock*}

  \begin{lstlisting}[basicstyle=\ttfamily\small]
[keystone:children]
control

[horizon:children]
control

[rabbitmq:children]
control

[mariadb:children]
database

...

[nova-compute:children]
compute
  \end{lstlisting}
\end{frame}

\begin{frame}[fragile]{Network control}
  \begin{columns}
    \begin{column}{.45\textwidth}
      \begin{lstlisting}
provider: g5k
resources:
  `\alert{grp1:}'
    paravance:
      control: 1
      database: 1
      nova-conductor: 1
      network: 1
      storage: 1
  `\alert{grp2:}'
    paravance:
      compute: 50
      \end{lstlisting}
    \end{column}
    \begin{column}{.45\textwidth}
      \begin{lstlisting}
`\alert{network\_constraints:}'
  `\alert{- src: grp1}'
    `\alert{dst: grp2}'
    `\alert{delay: 150ms}'
    `\alert{rate: 100mbit}'
    `\alert{symetric: true}'
      \end{lstlisting}
    \end{column}
  \end{columns}
\end{frame}

\begin{frame}[fragile]{OpenStack customization}
  \begin{lstlisting}
kolla:
  openstack_release: 4.0.0

patches:
  - name: patch mariadb bootstrap
    src: mariadb_bootstrap.yml
    dst: kolla/ansible/roles/mariadb/tasks/bootstrap.yml
    enabled: "yes"
  - name: patch galera.cnf.j2
    src: galera.cnf.j2
    dst: kolla/ansible/roles/mariadb/templates/galera.cnf.j2
    enabled: "yes"
  - name: patch haproxy.cfg.j2
    src: haproxy.cfg.j2
    dst: kolla/ansible/roles/haproxy/templates/haproxy.cfg.j2
    enabled: "yes"
  \end{lstlisting}
\end{frame}

\section{\texttt{enos bench}}

\begin{frame}[fragile]{\insertsectionhead}
  Execute arbitrary Rally benchmark
  \begin{lstlisting}
rally:
  enabled: true
  args:
    concurrency:
      - 5
      - 10
    times:
      - 100
    scenarios
      - name: boot and list servers
        file: nova-boot-list-cc.yml
  \end{lstlisting}
  \bigskip
  \texttt{\$ enos bench --workload=run.yml}
\end{frame}

%\begin{frame}{\insertsectionhead}
%  \textit{for each scenario}
%  \begin{itemize}
%	 \item Execute benchmark
%  \item Collect and store metrics
%    CPU/RAM/Network consumption whatever the service is (cAdvisor).
%  \item Generates reports
%  \end{itemize}
%\end{frame}

\section{\texttt{enos backup}}
\begin{frame}{\insertsectionhead}
  \texttt{enos backup} produces a archive with ready-to-share results:
  \begin{itemize}
  \item Rally/Shaker reports
  \item OpenStack logs
  \item InfluxDB database filled with collected metrics
  \item Grafana server for visualisation
  \end{itemize}
\end{frame}

\begin{frame}{Measures Output}
  \begin{center}
  \includegraphics[width=\textwidth]{figs/grafana}
  \end{center}
\end{frame}

%% \section{\texttt{enos info}}
%\begin{frame}[fragile]{Interact with Enos}
%  Enos fully exposes its state so you can add your own
%  tasks/benchmarks.
%  \bigskip
%  {
%    \scriptsize
%  \begin{lstlisting}
%$ enos info
%
%{
%  'cinder_database_password': 'demo',
%  ...
%  'mariadb_password': 'demo',
%  ...
%  'vip': '10.24.61.254',
%  'eths': ('eth0', 'eth1'),
%  'provider_net': {'cidr': '10.24.0.0/18',
%                   'start': '10.24.0.0',
%                   'gateway': ...},
%  'rsc': {'compute':[Host('paravance-16-kavlan-4.rennes.grid5000.fr'),
%                     Host('paravance-15-kavlan-4.rennes.grid5000.fr')],
%          'nova-conductor': ...}
%}
%  \end{lstlisting}
%  }
%\end{frame}

\begin{frame}{Conclusion}
  \textbf{Enos}
  \begin{itemize}
	 \item Complete solution to deploy and evaluate OpenStack
		\url{https://github.com/BeyondTheClouds/enos}
	 \item Integrated, reproducible and sharable
	 \item Fine control over deployment topology and network
	 \item Ready-to-use recipes:
		\url{https://github.com/BeyondTheClouds/enos-scenarios}
  \end{itemize}

  \vspace{.4cm}
  \textbf{Future works}
  \begin{itemize}
  \item Multiple Regions
  \item WANwide OpenStack
  \item OSProfiler Integration
  \item Adapt to other complex software stacks
  \end{itemize}
\end{frame}

\begin{frame}{Thank you}
  \begin{textblock*}{\paperwidth}(0cm, 1.0cm)
   \centering\includegraphics[scale=.7]{figs/discovery-sticker.png}
  \end{textblock*}

  \begin{textblock*}{.7\paperwidth}(1.5cm, 6.8cm)
  \begin{columns}
    \begin{column}{.65\textwidth}
		Ronan-Alexandre.Cherrueau@inria.fr\\
		Adrien.Lebre@inria.fr\\
		Dimitri.Pertin@inria.fr
    \end{column}
    \begin{column}{.55\textwidth}
		\underline{Anthony.Simonet@inria.fr}\\
		Matthieu.Simonin@inria.fr
    \end{column}
  \end{columns}
  \end{textblock*}
\end{frame}

\end{document}
